\section{Conclusion}
\label{sec:conclu}
% JHQ's Cartesian codebook makes its distance table exactly factorisable, and
% the measurements say where that pays. $G{=}2$ has the highest throughput in 23
% of 24 cells over six datasets, while $G{=}4$ and $G{=}8$ hold the same
% footprint but require different numbers of lookups --- so table size does not
% determine scan speed. Refinement accounts for $57\%$ and $70\%$ of measured
% work on Vogue and OpenAI-3072 at probe 32 and $\alpha=100$, motivating budget
% selection. These cross-dataset observations do not isolate a dimension effect.

% Sampled output agreement yields mean selected-arm QPS of
% $1.05$--$2.77\times$ the fixed-$\alpha$ arm in ten of twelve held-out cells. On arxiv it selects
% larger budgets that improve mean recall but lower throughput relative to
% fixed $\alpha=100$. Throughput gains do not establish a quality guarantee:
% at $S=128$, the fraction of draws losing more than $10^{-3}$ recall ranges
% from zero to 40/64. The evidence covers one GPU, $k=10$ and stationary query
% pools; drift, other $k$ and other devices are untested.




JHQ-GPU combines exact Cartesian table factorization with workload-level refinement calibration. Two-group tables lead in 23/24 configurations, while equal-footprint alternatives confirm that lookup work also matters. Calibration yields $1.05$--$2.77\times$ fixed-budget throughput in 10/12 held-out configurations; arxiv instead gains mean recall at lower throughput. Output agreement does not guarantee quality: at $S=128$, up to 40/64 draws lose more than $10^{-3}$ recall. Evidence is limited to one GPU, $k=10$, and stationary query pools; other devices, $k$, and workload drift remain untested.